\section{Projects}
\label{sec:projects}

A project file (\file{.raftcorr}) keeps a whole session: which images, the regions of interest, the
settings, the probes and, if you want, the results. Opening it brings the session back without
running the analysis again. The application saves and opens projects; \file{raftcorr run
-{}-project} also writes one, and the Python library reads and writes them
(Section~\ref{sec:cli}).

\subsection{Saving and opening}

\begin{description}[leftmargin=0.6cm,style=nextline]
  \item[\menu{File}{Save Project} (\kbd{Ctrl+S})]
    Saves to the project file of the session. The first time, it asks for a name, starting with
    \file{project.raftcorr} in the image folder; the extension \file{.raftcorr} is added. Images
    have to be loaded first.
  \item[\menu{File}{Save Project As…} (\kbd{Ctrl+Shift+S})]
    Always asks for a name, and the session continues in the new file.
  \item[\menu{File}{Open Project…} (\kbd{Ctrl+O})]
    Opens a project: the images, the regions of interest, the settings, the probes, the load
    record and the results it holds, on the frame that was shown when it was saved.
\end{description}

When the session has results, every save asks \ui{Include the computed results in this project?}
with the approximate uncompressed size of the results. \ui{Yes} (the default) stores the
displacement, the residual and the strain, so the project reopens ready to view, analyse and
export. \ui{No} writes a small file with the settings only, for sharing or as a template; after
opening it you run the analysis again. \ui{Cancel} saves nothing.

Saving and opening run behind a progress dialog. \ui{Cancel} in that dialog stops a save or an
open. A save is written to a temporary file first, so a cancelled or failed save leaves an
existing project as it was. An open reads the stored results frame by frame, the bar follows the
frames read, and a cancelled open leaves the current session as it was. After a save or an open, the title bar shows the name of the project file.

\subsection{What a project holds}

\begin{center}
\small
\begin{tabularx}{\linewidth}{@{}l X@{}}
  \toprule
  Stored & Content \\
  \midrule
  Images & the image folder (its full path and its position relative to the project file) and
    the file names in frame order; not the images themselves \\
  Regions of interest & the region of interest of frame~1 and of every later frame that has one \\
  Settings & tracking mode, model, iterations, smoothing, device, correlation and tiling settings;
    the strain settings; the physical units \\
  View & the displayed field, \ui{Show on}, background, opacity, \ui{Hide untracked regions}, each
    field's colormap and colour range, the current frame, and the last model suggestion \\
  Analysis & the probes with their names, colours and visibility; the load record of a testing
    machine (Section~\ref{sec:analysis}) \\
  Results (optional) & the displacement of every frame, the residual (in half precision), the run
    information, and the strain with its settings and trimmed band \\
  \bottomrule
\end{tabularx}
\end{center}

The images stay where they are. Keep the image folder with the project, or move both together.

\subsection{When the images have moved}

On opening, RAFTcorr looks for the image folder in this order:

\begin{enumerate}
  \item the folder saved in the project;
  \item the folder at the same position relative to the project file, for example after you moved
    or copied the project and the images together, or synced them to another computer. The
    \ui{LOG} says which folder was used;
  \item otherwise it says \ui{The image folder of this project was not found. Choose its new
    location.} and asks for the folder.
\end{enumerate}

The folder must hold every image file of the project under its saved name; if one is missing,
the project does not open and the message names the first missing file. A project saved by a
newer RAFTcorr, with a newer project format, does not open in an older one: the message asks you
to update RAFTcorr.

\subsection{Opening a project by double-click}
\label{sec:projects-doubleclick}

The Windows installer registers the \file{.raftcorr} extension, so a double-click on a project
file starts RAFTcorr and opens it. After a \file{pip} or source installation on Windows, use
\menu{File}{Associate .raftcorr files with RAFTcorr…} once. It registers the extension for your
user account only, needs no administrator rights, and confirms with \emph{``Done.
Double-clicking a .raftcorr file will now open RAFTcorr and restore that project.''}; if the
registry cannot be written, the dialog \ui{File Association Failed} gives the reason. Run it again
if you move or reinstall RAFTcorr, so that the double-click starts the copy you use. The menu
entry exists on Windows only.

On every platform, a project or an image folder can be given on the command line:

\begin{lstlisting}
raftcorr gui path/to/specimen.raftcorr
raftcorr gui path/to/images
\end{lstlisting}

\begin{tipbox}
Save with results before you close RAFTcorr after a long run. A project with results opens
without tracking anything; a settings-only project, or no project, means running the whole
sequence again.
\end{tipbox}
